% ============================================================
% DRAFT ONLY — do not merge automatically. Read this whole header,
% then copy the REPLACE-WITH text into introduction.tex by hand.
%
% WHAT THIS IS FOR
% I read the slide deck "ECSA - New GAs only.pdf" (the one you put in
% "Skripsie/New GA's/"). It's a lecture deck from Prof Arnold Rix on
% ECSA's revised Graduate Attribute (GA) definitions, explicitly flagged
% as relevant to PE448 for GA 1-6, 8 and 9 (GA7, 10, 11 are not required
% for PE448, same as before). The numbering of the GAs you already know
% about hasn't changed — it's the WORDING/scope of each one that has.
% The two shifts that actually land on your project, out of all nine:
%
% - GA5 (Use of engineering tools) used to be about using tools
%   competently. The new wording explicitly adds "recognise limitations
%   of appropriate techniques, resources and modern engineering and IT
%   tools" — i.e. an engineer is now expected to know and show WHERE a
%   tool breaks down, not just how to drive it. That is, almost word for
%   word, what RQ1/RQ2 and your whole evaluation chapter already do:
%   you're not just using a locally hosted LLM, you're mapping out
%   exactly where its planning capability runs out (depth, context,
%   lighting, etc).
% - GA9 (Independent learning) now explicitly calls out "adaptability to
%   new and emerging technologies" and research-literature engagement.
%   Locally hosted LLM/VLM robotics is about as "emerging" as it gets,
%   and you built the whole pipeline (ROS2 node graph reverse-engineered
%   from no docs, Ollama serving, the depth sweep methodology) largely
%   through your own independent reading rather than a taught module.
%
% I picked these two because they're the ones with a genuine, specific,
% non-generic match to what your project actually does — I didn't want
% to force-fit all seven required GAs into two sentences in the
% Introduction (that's what Appendix B is for, and I've already drafted
% that separately). If you'd rather anchor to different GAs here, tell
% me which and I'll redo it.
%
% WHERE IT GOES: end of \mysection{Background}, right before
% \mysection{Problem Statement and Research Questions} begins.
% ============================================================

% --- EDIT 1: insert new closing paragraph at the end of Background ---
%
% CURRENT (last paragraph of \mysection{Background}, ends with):
%
% Second, keeping inference off the rover and on a host machine that can run the same
% planner locally or swap to a cloud API, is what makes the local-versus-cloud comparison closer to
% apples-to-apples of model capability alone. If the rover itself hosted the model,
% switching the benchmark to a cloud API would also change the hardware story and the not a fair comparison.
% Third, it keeps the result general: the rover's own on-board computer, a Raspberry
% Pi, only has to stream camera frames and relay drive and sensor commands over the local network.
% Using the same rover would work unchanged with a different host machine entirely.
%
% \mysection{Problem Statement and Research Questions}{Problem Statement and Research Questions}
% \label{sec:problem-statement}
%
% REPLACE WITH (same text, plus one new paragraph inserted before the
% \mysection{Problem Statement...} line):
%
% Second, keeping inference off the rover and on a host machine that can run the same
% planner locally or swap to a cloud API, is what makes the local-versus-cloud comparison closer to
% apples-to-apples of model capability alone. If the rover itself hosted the model,
% switching the benchmark to a cloud API would also change the hardware story and the not a fair comparison.
% Third, it keeps the result general: the rover's own on-board computer, a Raspberry
% Pi, only has to stream camera frames and relay drive and sensor commands over the local network.
% Using the same rover would work unchanged with a different host machine entirely.
%
% This framing also lines up with ECSA's revised Graduate Attribute definitions for PE448: Graduate
% Attribute 5 (Use of Engineering Tools) now explicitly requires an engineer to recognise the
% limitations of the modern engineering and IT tools they apply, rather than only how to operate them,
% which is exactly what RQ1 and RQ2 set out to do for a locally hosted LLM and VLM pipeline. Building
% and evaluating that pipeline from scratch, with no existing local reference implementation to follow,
% also draws directly on Graduate Attribute 9 (Independent Learning Ability), in its updated emphasis
% on adapting to new and emerging technologies.
%
% \mysection{Problem Statement and Research Questions}{Problem Statement and Research Questions}
% \label{sec:problem-statement}
%
% ============================================================
% That's the only change in this file. Brace count below is for the
% inserted paragraph alone (it's plain prose, no macros that open a
% brace without closing it on the same line, so this should be 0/0 —
% checked anyway against the convention).
% ============================================================
